\paragraph{One field readout in the classical and quantum descriptions.}
The reconstruction of the charged action also fixes a class of observables
on its neutral Hilbert space. This identification uses the same dressed
field, volume and magnetic reconstruction in both descriptions; it does
not identify a classical trajectory with a quantum state.

Use the configuration space \(\mathcal Q\), metric volume \(d\mu\)
and neutral Hilbert space \(\mathcal H_0\) of
Theorem~\ref{thm:whitney-interacting-quantum}. For real
\(f\in L^\infty(\Omega)\) and a real square-integrable two-form
\(b\), define the configuration functions
\begin{equation}
 \begin{aligned}
 I_f(a,\psi)&=\int_\Omega f|\Psi_a|^2\,dx,\\
 J_f(a,\psi)&=\int_\Omega f|\Psi_a|^4\,dx,\\
 B_b(a,\psi)&=\int_\Omega\langle b,d(W_1a)\rangle\,dx.
 \end{aligned}
 \label{eq:whitney-common-observables}
\end{equation}
The form pairing and volume use the supplied cone metric. These functions
are finite and real on every configuration and are invariant under the
nodal gauge transformations. In particular, they are invariant under the
residual circle on \(\mathcal Q\). Spatial smearings specify which
field information is read; they do not supply a physical detector model.

\begin{theorem}[A common reconstructed observable algebra]
\label{thm:whitney-common-observables}
Let \(\mathcal O=(O_1,\ldots,O_k):\mathcal Q\to\mathbb R^k\)
be any finite list of functions from
\eqref{eq:whitney-common-observables}. Each multiplication operator
\[
 \widehat O_jF=O_jF,\qquad
 \mathcal D(\widehat O_j)
   =\{F\in\mathcal H_0:O_jF\in L^2(\mathcal Q,d\mu)\}
\]
is self-adjoint on \(\mathcal H_0\). They strongly commute, with
joint spectral projection-valued measure
\begin{equation}
 \mathsf P(B)F=\mathbf1_{\mathcal O^{-1}(B)}F,
 \qquad B\subseteq\mathbb R^k\text{ Borel}.
 \label{eq:whitney-common-observable-pvm}
\end{equation}
For every normalized \(F\in\mathcal H_0\), its spectral probability
law is
\begin{equation}
 \nu_F(B)=\langle F,\mathsf P(B)F\rangle
         =\int_{\mathcal O^{-1}(B)}|F(q)|^2\,d\mu(q).
 \label{eq:whitney-common-observable-law}
\end{equation}
For every bounded real Borel detector function \(d:\mathbb R^k\to
\mathbb R\), the operator \(d(\widehat{\mathcal O})\) is
multiplication by \(d\circ\mathcal O\). Thus the classical readout
\(d(\mathcal O(q))\) and its quantum spectral law use exactly the same
configuration function. No ground-state assumption is required.
\end{theorem}
\begin{proof}
The dressed interpolation is smooth in its finite coefficient variables,
with unit-modulus phases, and the cone has finite volume. On compact
coefficient sets its scalar fields and their coefficient derivatives are
uniformly bounded. Dominated convergence proves continuity of \(I_f,J_f\);
\(B_b\) is a continuous linear function of \(a\).
Gauge covariance of \(\Psi_a\) and \(d^2=0\) prove their invariance.

For any one such real function \(O\), truncation by
\(\mathbf1_{\{|O|\leq n\}}\) approximates every neutral Hilbert vector
and lies in the multiplication domain. This proves density. The operator
is symmetric. Multiplication by \((O\pm i)^{-1}\) is bounded and
preserves neutrality; for every \(F\in\mathcal H_0\) it gives a
domain vector mapped to \(F\) by \(\widehat O\pm i\).
Both resolvents are therefore everywhere defined, proving
self-adjointness. The same argument on \(\mathcal H\) gives its full
space realization. Invariance of \(O\) makes every spectral projection
commute with the residual-circle action and its neutral projection, so
\(\mathcal H_0\) is a reducing subspace.

The indicator multipliers in
\eqref{eq:whitney-common-observable-pvm} are orthogonal projections.
Disjoint countable unions give strong countable additivity by dominated
convergence, and \(\mathsf P(\mathbb R^k)=I_{\mathcal H_0}\).
Its coordinate spectral integrals have precisely the stated multiplication
domains. This proves joint spectrality and strong commutation. Integrating
a bounded Borel function proves the detector formula and
\eqref{eq:whitney-common-observable-law}.
\end{proof}

For any initial normalized neutral state, the established unitary group
gives \(F_t=e^{-it\widehat H_0/\hbar}F\), and the same law defines
\(\nu_{F_t}\) at every model time. This defines time-dependent
probabilities without claiming that those probabilities have been
numerically computed. For the Gaussian half-density in
Proposition~\ref{prop:whitney-interacting-gaussian-state}, all polynomial
coefficient moments are finite. In particular, the two exact scalar
moments in \eqref{eq:whitney-interacting-gaussian-matter-moments} are
the spectral expectations of \(\widehat I_1\) and \(\widehat J_1\)
from this construction.

The charged field \(\Psi_a\) itself is not circle invariant; its
multiplication operator does not in general preserve \(\mathcal H_0\).
Consequently the theorem concerns neutral field readouts, rather than
asserting that every gauge-fixed coefficient is a physical observable.
It also supplies no momentum-ordering prescription for a local quantum
current or kinetic-energy density. The spectral probability rule is that
of the declared Hilbert realization; its identification with experimental
outcomes remains a physical measurement assumption.

\begin{proposition}[Continuum convergence of neutral detector readouts]
\label{cor:whitney-continuum-detector-readouts}
Under the hypotheses and initialization of
Theorem~\ref{thm:whitney-real-continuum-trajectory}, fix finitely many real
bounded spatial smearings. The scalar readouts
\[
 I_f(u_n)=\int_\Omega f u_n^2\,dx,\qquad
 J_f(u_n)=\int_\Omega f u_n^4\,dx
\]
converge to the corresponding continuum readouts uniformly on \([0,T]\)
at rate \(O(n^{-1})\). Every fixed Lipschitz function of this finite
readout vector has the same rate. In particular, bounded Lipschitz
detector responses have a controlled classical continuum limit.
\end{proposition}
\begin{proof}
The trajectory theorem bounds \(u_n,u\) uniformly in \(H^1\), and
their difference there is \(O(n^{-1})\). Hence
\[
 |I_f(u_n)-I_f(u)|
 \leq\|f\|_\infty(\|u_n\|_2+\|u\|_2)\|u_n-u\|_2.
\]
Factorization of the fourth powers and H\"older's inequality give
\[
 |J_f(u_n)-J_f(u)|
 \leq C\|f\|_\infty(\|u_n\|_4^3+\|u\|_4^3)
                                      \|u_n-u\|_4.
\]
Use the fixed-domain embedding \(H^1\hookrightarrow L^4\) and then
the Lipschitz inequality for the detector function.
\end{proof}

The continuum corollary concerns the full action's invariant real sector,
where electromagnetic fields and current vanish. It proves no convergence
of the quantum states or spectral measures across refinements. No
quantum--classical state correspondence, calibrated detector, physical
spacetime identification or source-selected action follows from sharing
the observable functions. The common input is the stated cone action and
its reconstruction and quantization rules.
